% GENERATED FILE. Edit canonical lessons, then run npm run generate:books.
% canonical-source-hash: fnv1a64:596a4077c9f946da
% canonical-lessons: ML-C464-kasrramsu, ML-C464-emigresan, ML-C464-cekkposrru, ML-C464-parisodhana, ML-C464-videsi

\chapter{At the Border}
\label{ch:ml-at-the-border}
\hlchaptermodality{464}

\begin{chapteropening}
\textbf{By the end of this chapter:} \emph{I can say customs, passport control, a checkpost, an inspection and a foreigner.}

Complete the last lesson of chapter 464: I can say customs, passport control, a checkpost, an inspection and a foreigner.
\end{chapteropening}

\section[\texorpdfstring{\ml{കസ്റ്റംസ്} (kasṟṟaṁsŭ)}{കസ്റ്റംസ് (kasṟṟaṁsŭ)}]{\ml{കസ്റ്റംസ്} (kasṟṟaṁsŭ) — customs}

\label{lesson:ML-C464-kasrramsu}



\begin{quote}
Before the new one: say the Malayalam for to lift off, then the Malayalam for a luggage trolley.
\end{quote}

\subsection*{You'll want to know: \ml{കസ്റ്റംസ്}}
\textbf{\ml{കസ്റ്റംസ്}} — \emph{kasṟṟaṁsŭ} — \textquotedblleft{}customs\textquotedblright{}.

\begin{quote}
\ml{കസ്റ്റംസ്} \ml{ഉദ്യോഗസ്ഥൻ} \ml{എന്റെ} \ml{ബാഗ്} \ml{തുറന്നു}.
\end{quote}

\begin{quote}
\emph{kasṟṟaṁsŭ udyōgasthan enṟe bāgŭ tuṟannu.}
\end{quote}

\begin{quote}
The customs officer opened my bag.
\end{quote}

A customs officer is a \ml{കസ്റ്റംസ്} \ml{ഉദ്യോഗസ്ഥൻ}.

\begin{cousinweb}
\textbf{\ml{കസ്റ്റംസ്}} (\emph{kasṟṟaṁsŭ}) was borrowed from English \emph{customs}, the office that checks goods at a border.
\end{cousinweb}

\begin{grammarlens}[title={What you've built}]
One more word for getting around.
\end{grammarlens}

\subsection*{Your turn}
\begin{itemize}
\raggedright
  \item \emph{Say it:} \emph{kasṟṟaṁsŭ}
  \item \emph{Say it:} \emph{kasṟṟaṁsŭ}, once more
  \item \emph{Say it:} \emph{kasṟṟaṁsŭ udyōgasthan enṟe bāgŭ tuṟannu.} — The customs officer opened my bag.
  \item \emph{Recall:} say the Malayalam for to lift off, then the Malayalam for a luggage trolley, then say \emph{kasṟṟaṁsŭ} again
\end{itemize}

\begin{tcolorbox}[breakable,colback=teal!4,colframe=teal!35!black,title={Before you move on}]
What does \ml{കസ്റ്റംസ്} mean? (\textquotedblleft{}Customs\textquotedblright{}.) Say it once more.
\end{tcolorbox}
\section[\texorpdfstring{\ml{എമിഗ്രേഷൻ} (emigrēṣan)}{എമിഗ്രേഷൻ (emigrēṣan)}]{\ml{എമിഗ്രേഷൻ} (emigrēṣan) — passport control}

\label{lesson:ML-C464-emigresan}



\begin{quote}
Before the new one: say the Malayalam for a luggage trolley, then the Malayalam for customs.
\end{quote}

\subsection*{You'll want to know: \ml{എമിഗ്രേഷൻ}}
\textbf{\ml{എമിഗ്രേഷൻ}} — \emph{emigrēṣan} — \textquotedblleft{}passport control\textquotedblright{}.

\begin{quote}
\ml{എമിഗ്രേഷനിൽ} \ml{അവർ} \ml{എന്റെ} \ml{പാസ്പോർട്ടും} \ml{വിസയും} \ml{പരിശോധിച്ചു}.
\end{quote}

\begin{quote}
\emph{emigrēṣanil avar enṟe pāspōrṭṭuṁ visayuṁ pariśōdhiccu.}
\end{quote}

\begin{quote}
At immigration they checked my passport and visa.
\end{quote}

Malayalam speakers usually call the passport desk at an airport \ml{എമിഗ്രേഷൻ}, the English name of the check on people leaving.

\begin{cousinweb}
\textbf{\ml{എമിഗ്രേഷൻ}} (\emph{emigrēṣan}) was borrowed from English \emph{emigration}, the passport check for people leaving a country.
\end{cousinweb}

\begin{grammarlens}[title={What you've built}]
One more word for getting around.
\end{grammarlens}

\subsection*{Your turn}
\begin{itemize}
\raggedright
  \item \emph{Say it:} \emph{emigrēṣan}
  \item \emph{Say it:} \emph{emigrēṣan}, once more
  \item \emph{Say it:} \emph{emigrēṣanil avar enṟe pāspōrṭṭuṁ visayuṁ pariśōdhiccu.} — At immigration they checked my passport and visa.
  \item \emph{Recall:} say the Malayalam for a luggage trolley, then the Malayalam for customs, then say \emph{emigrēṣan} again
\end{itemize}

\begin{tcolorbox}[breakable,colback=teal!4,colframe=teal!35!black,title={Before you move on}]
What does \ml{എമിഗ്രേഷൻ} mean? (\textquotedblleft{}Passport control\textquotedblright{}.) Say it once more.
\end{tcolorbox}
\section[\texorpdfstring{\ml{ചെക്ക്പോസ്റ്റ്} (cekkpōsṟṟŭ)}{ചെക്ക്പോസ്റ്റ് (cekkpōsṟṟŭ)}]{\ml{ചെക്ക്പോസ്റ്റ്} (cekkpōsṟṟŭ) — a checkpost}

\label{lesson:ML-C464-cekkposrru}



\begin{quote}
Before the new one: say the Malayalam for customs, then the Malayalam for passport control.
\end{quote}

\subsection*{You'll want to know: \ml{ചെക്ക്പോസ്റ്റ്}}
\textbf{\ml{ചെക്ക്പോസ്റ്റ്}} — \emph{cekkpōsṟṟŭ} — \textquotedblleft{}a checkpost\textquotedblright{}.

\begin{quote}
\ml{അതിർത്തിയിലെ} \ml{ചെക്ക്പോസ്റ്റിൽ} \ml{ബസ്} \ml{അര} \ml{മണിക്കൂർ} \ml{നിന്നു}.
\end{quote}

\begin{quote}
\emph{atirttiyile cekkpōsṟṟil basŭ ara maṇikkūr ninnu.}
\end{quote}

\begin{quote}
The bus stood for half an hour at the checkpost on the border.
\end{quote}

A checkpost stands on a road where it crosses from one state or district into another.

\begin{cousinweb}
\textbf{\ml{ചെക്ക്പോസ്റ്റ്}} (\emph{cekkpōsṟṟŭ}) was borrowed from English \emph{check-post}, a post where vehicles are stopped and checked.
\end{cousinweb}

\begin{grammarlens}[title={What you've built}]
One more word for getting around.
\end{grammarlens}

\subsection*{Your turn}
\begin{itemize}
\raggedright
  \item \emph{Say it:} \emph{cekkpōsṟṟŭ}
  \item \emph{Say it:} \emph{cekkpōsṟṟŭ}, once more
  \item \emph{Say it:} \emph{atirttiyile cekkpōsṟṟil basŭ ara maṇikkūr ninnu.} — The bus stood for half an hour at the checkpost on the border.
  \item \emph{Recall:} say the Malayalam for customs, then the Malayalam for passport control, then say \emph{cekkpōsṟṟŭ} again
\end{itemize}

\begin{tcolorbox}[breakable,colback=teal!4,colframe=teal!35!black,title={Before you move on}]
What does \ml{ചെക്ക്പോസ്റ്റ്} mean? (\textquotedblleft{}A checkpost\textquotedblright{}.) Say it once more.
\end{tcolorbox}
\section[\texorpdfstring{\ml{പരിശോധന} (pariśōdhana)}{പരിശോധന (pariśōdhana)}]{\ml{പരിശോധന} (pariśōdhana) — an inspection; a check}

\label{lesson:ML-C464-parisodhana}



\begin{quote}
Before the new one: say the Malayalam for passport control, then the Malayalam for a checkpost.
\end{quote}

\subsection*{You'll want to know: \ml{പരിശോധന}}
\textbf{\ml{പരിശോധന}} — \emph{pariśōdhana} — \textquotedblleft{}an inspection; a check\textquotedblright{}.

\begin{quote}
\ml{വിമാനത്താവളത്തിൽ} \ml{സുരക്ഷാ} \ml{പരിശോധന} \ml{ഉണ്ടായിരുന്നു}.
\end{quote}

\begin{quote}
\emph{vimānattāvaḷattil surakṣā pariśōdhana uṇṭāyirunnu.}
\end{quote}

\begin{quote}
There was a security check at the airport.
\end{quote}

The noun of \ml{പരിശോധിക്കുക}, to check. A medical check-up is also a \ml{പരിശോധന}.

\begin{cousinweb}
\textbf{\ml{പരിശോധന}} (\emph{pariśōdhana}) is borrowed from Sanskrit \emph{pariśodhana}, \textquotedblleft{}cleaning, purification\textquotedblright{} (Monier-Williams). In Malayalam it means a close check or an inspection.
\end{cousinweb}

\begin{grammarlens}[title={What you've built}]
One more word for getting around.
\end{grammarlens}

\subsection*{Your turn}
\begin{itemize}
\raggedright
  \item \emph{Say it:} \emph{pariśōdhana}
  \item \emph{Say it:} \emph{pariśōdhana}, once more
  \item \emph{Say it:} \emph{vimānattāvaḷattil surakṣā pariśōdhana uṇṭāyirunnu.} — There was a security check at the airport.
  \item \emph{Recall:} say the Malayalam for passport control, then the Malayalam for a checkpost, then say \emph{pariśōdhana} again
\end{itemize}

\begin{tcolorbox}[breakable,colback=teal!4,colframe=teal!35!black,title={Before you move on}]
What does \ml{പരിശോധന} mean? (\textquotedblleft{}An inspection; a check\textquotedblright{}.) Say it once more.
\end{tcolorbox}
\section[\texorpdfstring{\ml{വിദേശി} (vidēśi)}{വിദേശി (vidēśi)}]{\ml{വിദേശി} (vidēśi) — a foreigner}

\label{lesson:ML-C464-videsi}



\begin{quote}
Before the new one: say the Malayalam for a checkpost, then the Malayalam for an inspection.
\end{quote}

\subsection*{You'll want to know: \ml{വിദേശി}}
\textbf{\ml{വിദേശി}} — \emph{vidēśi} — \textquotedblleft{}a foreigner\textquotedblright{}.

\begin{quote}
\ml{ആ} \ml{ഹോട്ടലിൽ} \ml{ധാരാളം} \ml{വിദേശികൾ} \ml{താമസിക്കുന്നു}.
\end{quote}

\begin{quote}
\emph{ā hōṭṭalil dhārāḷaṁ vidēśikaḷ tāmasikkunnu.}
\end{quote}

\begin{quote}
Many foreigners stay at that hotel.
\end{quote}

Plural: \ml{വിദേശികൾ}, \emph{vidēśikaḷ}.

Built on \ml{വിദേശം}, abroad: the -\ml{ഇ} ending names a person from there.

\begin{cousinweb}
\textbf{\ml{വിദേശി}} (\emph{vidēśi}) is built from Sanskrit \emph{videśa}, \textquotedblleft{}another country, abroad\textquotedblright{} (Monier-Williams), with the ending -\ml{ഇ} for a person of that place.
\end{cousinweb}

\begin{grammarlens}[title={What you've built}]
One more word for getting around.
\end{grammarlens}

\subsection*{Your turn}
\begin{itemize}
\raggedright
  \item \emph{Say it:} \emph{vidēśi}
  \item \emph{Say it:} \emph{vidēśi}, once more
  \item \emph{Say it:} \emph{ā hōṭṭalil dhārāḷaṁ vidēśikaḷ tāmasikkunnu.} — Many foreigners stay at that hotel.
  \item \emph{Recall:} say the Malayalam for a checkpost, then the Malayalam for an inspection, then say \emph{vidēśi} again
\end{itemize}

\begin{tcolorbox}[breakable,colback=teal!4,colframe=teal!35!black,title={Before you move on}]
What does \ml{വിദേശി} mean? (\textquotedblleft{}A foreigner\textquotedblright{}.) Say it once more.
\end{tcolorbox}
